\documentclass[12pt,a4paper]{article}
\usepackage[utf8]{inputenc}
\usepackage[russian]{babel}
\usepackage{amsmath,amssymb,amsthm}
\usepackage{geometry}
\usepackage{graphicx}
\geometry{left=2cm,right=2cm,top=2cm,bottom=2cm}

\title{Субдифференциал и субградиент}
\author{}
\date{}

\begin{document}

\maketitle

\textbf{Задача 1. (2 балла)} Найдите субдифференциалы функций $f: \mathbb{R} \to \mathbb{R}$:

\textbf{1.1. (1 балл)}

\[
f(x) = \max \{-x, x, x^2\}.
\]

\textbf{1.1 решение:}

график функции выглядит таким образом

\vspace{1em}

\includegraphics[width=0.3\linewidth]{1.1.png}

\vspace{1em}

Понятно что субдифференциал в дифференцируемых функциях будет просто содержать только производную в этой точке.

В не дифференцируемых точках субдифф будет содержать множество наклонов удовлетворяющие условию. Например в точке $x = -1 \ \rightarrow$ $x^2' < x'$ поэтому $\frac{\partial f}{\partial x} = (x^2'; x')$


$$
\frac{\partial f}{\partial x} = 
\begin{cases}
  \{2x\}, & x \in (-\infty ; -1) \cup(1; + \infty) \\
  [-2; -1],  & x = -1 \\
  {-1}, & x \in (-1; 0) \\
  [-1; 1], & x = 0 \\
  {1}, & x \in (0; 1) \\
  [1; 2], & x = 1
\end{cases}
$$


\textbf{1.2. (1 балл)}

\[
f(x) = |x-2| + |x+2| + |x-1| + x^2.
\]

\textbf{1.2 решение:}

$$
|x-2| + |x+2| + |x-1| + x^2 = 
\begin{cases}
  x^2 - 3x + 1, & x < -2 \\
  x^2 - x + 5,  & x \in [-2; 1] \\
  x^2 + x + 3, & x \in [1; 2] \\
  x^2 + 3x - 1, & x > 2 
\end{cases}
$$

$$
\nabla f(x) = 
\begin{cases}
  2x -3, & x < -2 \\
  2x -1,  & x \in [-2; 1] \\
  2x + 1, & x \in [1; 2] \\
  2x + 3, & x > 2 
\end{cases}
$$

тогда 

$$
\frac{\partial f}{\partial x} = 
\begin{cases}
  2x -3, & x < -2 \\
  2x -1,  & x \in [-2; 1] \\
  2x + 1, & x \in [1; 2] \\
  2x + 3, & x > 2 \\
  [-7; -5], & x = -2 \\
  [1; 3], & x = 1 \\
  [5; 7], & x = 2
\end{cases}
$$

\textbf{Задача 2. (2 балла)} Найдите субдифференциал функции $f: \mathbb{R}^2 \to \mathbb{R}$

\[
f(x, y) = x^2 + xy + y^2 + 3|x + y - 2|.
\]

Также найдите точки минимума функции, используя полученный результат.


\textbf{2 Решение}

$$
f(x, y) = g(x, y) + h(x, y), \quad
g(x, y) = x^2 + xy + y^2, \quad
h(x, y) = 3|x + y - 2|
$$

$$
\partial f(x, y) = \nabla g(x, y) + \partial h(x, y)
$$

$$
\nabla g(x, y) = (2x + y, x + 2y)
$$

$$
\frac{\partial h}{\partial x} = 
\begin{cases}
    {-1}, & x + y < 2 \\
    [-1, 1], & x+y = 2 \\
    {1}, & x +y > 2
    \end{cases}
$$

$$
\frac{\partial h}{\partial y} = 
\begin{cases}
    {-1}, & x + y < 2 \\
    [-1, 1], & x + y = 2 \\
    {1}, & x +y > 2
    \end{cases}
$$

$$
\frac{\partial f}{\partial (x, y)} = 
\begin{cases}
    \{(2x + y - 3, x + 2y - 3)\}, & x + y < 2 \\
    \{(2x + y + 3 \lambda, x + 2y + 2 \lambda) \lambda \in [-1; 1]\}, & x + y = 2 \\
    \{(2x + y + 3, x + 2y + 3)\}, & x + y > 2
\end{cases}
$$

\vspace{1em}

во всех случаях кроме $x + y = 2$, субградиент не равен 0, рассмотрим случай $x + y = 2$

$$
\begin{cases}
2x + y + 3\lambda = 0 \\
x + 2y + 3\lambda = 0 \\
x + y = 2
\end{cases}
$$

отсюда и получаем что $x = y = 1, \quad \lambda = -1$

Значит $(x, y) = (1, 1)$ - минимум.

\textbf{Задача 3. (4 балла)} Найдите субдифференциалы функций $f: \mathbb{R}^d \to \mathbb{R}$:

\textbf{3.1. (1 балл)}

\[
f(x) = \log \left(1 + e^{\|x\|_1}\right).
\]

\textbf{3.1 Решение}

$$
\begin{aligned}
  &\partial f(x) = \\
  &\partial \log \left(1 + e^{x}\right) \cdot \partial \| x \|_1 =\\
  &\frac{e^{\|x\|_1}}{1 + e^{\|x\|_1}} \cdot \partial \| x \|_1
\end{aligned}
$$

где

$$
\partial \|x\|_1 = 
\begin{cases}
B_{\|\cdot\|_1^*}(0,1), & x = 0, \\
\{ g \in \mathbb{R}^d : \|g\|_1^* \leq 1,\ \langle g, x \rangle = \|x\|_1\}, & x \neq 0.
\end{cases}
$$

\textbf{3.2. (1 балл)}

\[
f(x) = \bigl(\max (0, \langle Ax, x\rangle + \langle b, x\rangle + c)\bigr)^3,
\]

где $A \in \mathbb{S}^d_{++}$, $b \in \mathbb{R}^d$, $c \in \mathbb{R}$.

\textbf{3.2 Решение}

пусть 
$$g_1(x) = 0, \quad g_2(x) = \langle Ax, x\rangle + \langle b, x\rangle + c$$

Тогда 
$$\partial g_1(x) = {0} \quad \partial g_2(x) = {2Ax + b}$$

в эллипсоиде $S = \{x : \langle Ax, x\rangle + \langle b, x\rangle + c = 0\}  \rightarrow g_1(x) = g_2(x)$

$$
g(x) = max(g_1(x), g_2(x))
$$

по теореме Дубовицкого-Милютина 

$$
\partial g(x) = 
\begin{cases}
    {0}, & \langle Ax, x\rangle + \langle b, x\rangle + c < 0 \\
    conv\{0, 2Ax + b\}, & \langle Ax, x\rangle + \langle b, x\rangle + c = 0 \\
    2Ax + b, & \langle Ax, x\rangle + \langle b, x\rangle + c > 0
\end{cases}
$$

$$
\partial f(x) = 3 \max (0, \langle Ax, x\rangle + \langle b, x\rangle + c))^2 * \partial g(x)
$$


\textbf{3.3. (2 балла)}

\[
f(x) = \left( \sum_{i=1}^d \max(0, x_i) \right)^2.
\]

\textbf{3.3 Решение}

Под суммой будет подразумеватся сумма Минковского:

$$
\begin{aligned}
  &\partial f(x) = \partial((\sum_{i=1}^d \max(0, x_i))^2) = \\
  &= 2 * (\sum_{i=1}^d \max(0, x_i)) \cdot \partial(\sum_{i=1}^d \max(0, x_i)) 
\end{aligned}
$$

$\partial(\sum_{i=1}^d \max(0, x_i)) = \{ \partial \max(0, x_i) \}$

$$
\partial \max(0, x_i)=
\begin{cases}
    \{0\}, & x_i < 0 \\
    [0, 1], & x_i = 0 \\
    \{1\}, & x_i > 0
\end{cases}
$$

\textbf{Задача 4. (2 балла)} Покажите, что субдифференциал в нуле функции $f: \mathbb{R}^d \to \mathbb{R}$:

\[
f(x) = \|x\|_\infty
\]

равен

\[
\partial f(0) = \operatorname{conv} \{ \pm e_1, \dots, \pm e_d \}.
\]


\textbf{4 Решение}

мы знаем как выглядит субдифференциал в нуле

$$
\partial \|x\|_1 = 
\begin{cases}
B_{\|\cdot\|_1^*}(0,1), & x = 0, \\
\{ g \in \mathbb{R}^d : \|g\|_1^* \leq 1,\ \langle g, x \rangle = \|x\|_1\}, & x \neq 0.
\end{cases}
$$

$\| g \|_1^*  = \sum_{i=1}^d |g_i| < 1$
Получается что g является линейной комбинацией $e_i$

\end{document}
